
\section{原始数据与公共物理模型核验}
\label{app:common}

\subsection{数据对象、单位与复现口径}

全文公共模型使用同一套节点、货箱、机型、DEM与任务参数。内部统一采用 m、s、kg、m$^3$、kWh 和 dB；正文显示时允许换算为 min，但后续计算始终使用未舍入原值。Q1公共物理验证不读取Q2--Q4求解结果，Q2--Q4均继承同一地形、航段与能耗口径。

\begin{table}[H]
\centering
\caption{公共模型关键核验结果}
\label{tab:app_common_validation}
\begin{tabular}{p{4.2cm}p{7.0cm}p{3.2cm}}
\toprule
核验对象 & 独立核验方式 & 结果\\
\midrule
TIF/MAT高程矩阵 & 逐像元比对 & 完全一致\\
AEQD水平距离 & GeographicLib独立测地线计算 & 最大差 $1.412\times10^{-9}$ m\\
DEM航段像元集合 & 网格边界交点遍历与条带区间求交 & 完全一致\\
45个安全载荷根 & Brent法与独立二分法 & 最大差 $1.843\times10^{-11}$ kg\\
Q1精确优化 & 原始箱级MILP、压缩MILP、DP & 词典序目标一致\\
自动化测试 & 42项测试 & 0失败、0错误、0跳过\\
\bottomrule
\end{tabular}
\end{table}

\subsection{公共物理错误注入}

为验证公共物理层检查器确实能够识别错误，而非仅重放既有结果，对以下七类错误进行了故障注入。

\begin{table}[H]
\centering
\caption{公共物理模型故障注入}
\label{tab:app_fault_injection}
\begin{tabular}{p{4.4cm}p{9.2cm}}
\toprule
错误类型 & 检出机制\\
\midrule
将载荷--航程指数 $3/2$ 错改为1 & 独立航程中间值与能耗重算不一致\\
遗漏返程爬升 & 往返能耗独立检查失败\\
返程错误沿用起飞载荷 & 逐段载荷重算失败\\
跨服务区混箱 & 服务区归属检查失败\\
重复分配同一货箱 & 原始货箱ID唯一计数失败\\
只检查质量、忽略体积 & 质量可行但体积超限的批型被拒绝\\
使用thin Bresenham替代完整像元遍历 & 航段相交像元集合不一致\\
\bottomrule
\end{tabular}
\end{table}

\subsection{复现性}

Q1在全新输出目录中从原始输入完整重跑，27个确定性数值文件与5幅生成图逐字节一致。求解耗时、平台元数据和可交换箱号代表不要求字节一致。公共模型的目标是保证“任务属性算得一致”，而各问最优性或近优性证据分别在后续附录中给出。

% ============================================================
\section{问题一完整能力、组批与敏感性结果}
\label{app:q1}

\subsection{45组最大安全载荷}

20\%返航安全余量下，15个服务区与3种机型的最大安全载荷如下。除表中低于额定载荷的6项外，其余组合均由额定载荷而非能源约束决定。

\begin{table}[H]
\centering
\caption{Q1各服务区三类机型最大安全载荷/kg}
\label{tab:app_q1_payload}
\begin{tabular}{lrrr}
\toprule
服务区 & A型 & B型 & C型\\
\midrule
S001 & 25.000000 & 30.000000 & 80.000000\\
S002 & 25.000000 & 30.000000 & 68.860343\\
S003 & 25.000000 & 30.000000 & 68.206491\\
S004 & 25.000000 & 30.000000 & 63.696955\\
S005 & 25.000000 & 30.000000 & 80.000000\\
S006 & 25.000000 & 30.000000 & 80.000000\\
S007 & 25.000000 & 30.000000 & 80.000000\\
S008 & 25.000000 & 28.801234 & 58.904361\\
S009 & 25.000000 & 30.000000 & 80.000000\\
S010 & 25.000000 & 30.000000 & 80.000000\\
S011 & 25.000000 & 30.000000 & 80.000000\\
S012 & 25.000000 & 30.000000 & 68.117193\\
S013 & 25.000000 & 30.000000 & 80.000000\\
S014 & 25.000000 & 30.000000 & 80.000000\\
S015 & 25.000000 & 30.000000 & 80.000000\\
\bottomrule
\end{tabular}
\end{table}

\subsection{18架次正式组批摘要}

完整箱号分配可由机器可读文件 \path{problems/D/q1/results/optimal_batches.csv} 复现。表\ref{tab:app_q1_batches}给出18个正式批次的主要物理属性。

\begin{longtable}{llllrrrr}
\caption{Q1正式18架次批次摘要}\label{tab:app_q1_batches}\\
\toprule
批次 & 区域 & 机型 & 质量/kg & 体积/m$^3$ & 能耗/kWh & 作业时间/s & 返航SOC/\%\\
\midrule
\endfirsthead
\toprule
批次 & 区域 & 机型 & 质量/kg & 体积/m$^3$ & 能耗/kWh & 作业时间/s & 返航SOC/\%\\
\midrule
\endhead
Q1-001 & S001 & C & 76 & 0.170 & 2.917582 & 1494.169 & 63.530\\
Q1-002 & S001 & C & 78 & 0.224 & 2.995769 & 1692.169 & 62.553\\
Q1-003 & S002 & B & 25 & 0.067 & 2.713422 & 1816.694 & 32.164\\
Q1-004 & S002 & C & 56 & 0.144 & 5.721225 & 2041.812 & 28.485\\
Q1-005 & S003 & B & 25 & 0.067 & 2.779917 & 2080.523 & 30.502\\
Q1-006 & S003 & C & 56 & 0.144 & 5.764965 & 2370.091 & 27.938\\
Q1-007 & S004 & C & 59 & 0.156 & 6.152932 & 2301.057 & 23.088\\
Q1-008 & S005 & C & 59 & 0.156 & 4.222193 & 1877.191 & 47.223\\
Q1-009 & S006 & C & 59 & 0.156 & 2.597988 & 1561.554 & 67.525\\
Q1-010 & S007 & C & 45 & 0.129 & 3.410894 & 1907.566 & 57.364\\
Q1-011 & S008 & C & 45 & 0.129 & 5.836012 & 2199.701 & 27.050\\
Q1-012 & S009 & B & 25 & 0.067 & 2.096234 & 1558.442 & 47.594\\
Q1-013 & S010 & B & 25 & 0.067 & 1.823196 & 1555.380 & 54.420\\
Q1-014 & S011 & B & 25 & 0.067 & 1.073952 & 1188.420 & 73.151\\
Q1-015 & S012 & B & 25 & 0.067 & 2.752414 & 1922.740 & 31.190\\
Q1-016 & S013 & B & 25 & 0.067 & 1.824759 & 1510.472 & 54.381\\
Q1-017 & S014 & B & 25 & 0.067 & 2.140388 & 1869.756 & 46.490\\
Q1-018 & S015 & B & 25 & 0.067 & 2.307453 & 1828.277 & 42.314\\
\bottomrule
\end{longtable}

\begin{figure}[H]
  \centering
  \includegraphics[width=0.82\textwidth]{map_single_routes.png}
  \caption{Q1基地O01至15个服务区的单点往返航线与DEM背景示意。该图仅用于空间关系展示，正式地形计算采用完整DEM像元相交。}
  \label{fig:app_q1_routes}
\end{figure}

\subsection{基线与Pareto结果}

\begin{table}[H]
\centering
\caption{Q1主要基线与正式方案}
\label{tab:app_q1_baseline}
\begin{tabular}{llrrr}
\toprule
机型策略 & 方法 & 架次 & 能耗/kWh & 累计作业时间/min\\
\midrule
A-only & exact & 52 & 112.180153 & 1488.763007\\
B-only & exact & 35 & 68.350657 & 927.625473\\
C-only & exact & 18 & 73.934908 & 563.979235\\
mixed & FFD-mass & 18 & 75.116697 & 563.979235\\
mixed & exact & 18 & \textbf{59.131296} & \textbf{546.266929}\\
\bottomrule
\end{tabular}
\end{table}

完整不同Pareto目标向量仅有两组：
\[
(18,\ 59.131296022053,\ 546.266929148593),
\]
\[
(19,\ 59.033942092203,\ 573.999623452513).
\]

\subsection{改变全局最少架次数的安全余量阈值}

\begin{table}[H]
\centering
\caption{Q1八个全局架次数临界阈值}
\label{tab:app_q1_thresholds}
\begin{tabular}{lcrr}
\toprule
安全余量阈值/\% & 触发区域 & 阈值处架次 & 严格右侧\\
\midrule
23.088350 & S004 & 18 & 19\\
27.049854 & S008 & 19 & 20\\
33.625809 & S008 & 20 & 21\\
33.896918 & S003 & 21 & 22\\
34.373532 & S004 & 22 & 23\\
34.391873 & S002 & 23 & 24\\
34.777897 & S008 & 24 & 25\\
35.267807 & S008 & 25 & 不可行\\
\bottomrule
\end{tabular}
\end{table}

完整569个候选事件与19次词典序最优代表变化由 \path{problems/D/q1/results/reserve_thresholds.csv} 和 \path{problems/D/q1/results/reserve_optimal_segments.csv} 保存。

% ============================================================
\section{问题二任务构造、调度结果与认证证据}
\label{app:q2}

\subsection{23个正式运输任务}

表\ref{tab:app_q2_tasks}直接对应 \path{q2_final/solution/main_23/decisions.json}。任务编号T01--T23仅用于论文展示，原始箱号保持不变。

\scriptsize
\begin{longtable}{cllp{8.2cm}}
\caption{Q2正式23个运输任务构成}\label{tab:app_q2_tasks}\\
\toprule
任务 & 机型 & 访问顺序 & 原始货箱ID\\
\midrule
\endfirsthead
\toprule
任务 & 机型 & 访问顺序 & 原始货箱ID\\
\midrule
\endhead
T01 & C & S001 & S001-MED-01, S001-MED-02, S001-WAT-06, S001-WAT-07, S001-WAT-08, S001-FOD-03, S001-HYG-01, S001-HYG-02\\
T02 & C & S001 & S001-WAT-02, S001-WAT-03, S001-WAT-04, S001-WAT-05, S001-FOD-02\\
T03 & B & S010 & S010-MED-01, S010-WAT-01, S010-FOD-01\\
T04 & B & S014 & S014-MED-01, S014-WAT-01, S014-FOD-01\\
T05 & B & S012 & S012-MED-01, S012-WAT-01, S012-FOD-01\\
T06 & A & S001 & S001-WAT-01, S001-FOD-01\\
T07 & A & S004 & S004-WAT-01, S004-FOD-01\\
T08 & A & S013 & S013-WAT-01, S013-FOD-01\\
T09 & A & S011$\rightarrow$S004 & S004-MED-01, S004-WAT-02, S011-MED-01\\
T10 & C & S006 & S006-MED-01, S006-WAT-01, S006-WAT-02, S006-WAT-03, S006-FOD-01, S006-HYG-01\\
T11 & C & S002 & S002-MED-01, S002-WAT-02, S002-WAT-03, S002-WAT-04, S002-FOD-01, S002-FOD-02, S002-HYG-01\\
T12 & C & S007$\rightarrow$S003 & S003-WAT-02, S003-WAT-03, S003-WAT-04, S003-FOD-01, S003-FOD-02, S003-HYG-01, S007-HYG-01\\
T13 & C & S005 & S005-MED-01, S005-WAT-01, S005-WAT-02, S005-WAT-03, S005-FOD-01, S005-HYG-01\\
T14 & A & S009$\rightarrow$S002$\rightarrow$S013 & S002-WAT-01, S009-MED-01, S013-MED-01\\
T15 & A & S007 & S007-MED-01, S007-WAT-01\\
T16 & A & S007 & S007-WAT-02, S007-FOD-01\\
T17 & B & S004 & S004-WAT-03, S004-HYG-01\\
T18 & B & S008 & S008-MED-01, S008-WAT-01, S008-FOD-01\\
T19 & B & S008 & S008-WAT-02, S008-HYG-01\\
T20 & A & S003$\rightarrow$S015 & S003-MED-01, S003-WAT-01, S015-MED-01\\
T21 & A & S015 & S015-WAT-01, S015-FOD-01\\
T22 & A & S009 & S009-WAT-01, S009-FOD-01\\
T23 & A & S011 & S011-WAT-01, S011-FOD-01\\
\bottomrule
\end{longtable}
\normalsize

\subsection{正式主方案与22架次见证}

\begin{table}[H]
\centering
\caption{Q2正式23架次方案与22架次可行见证}
\label{tab:app_q2_22_23}
\begin{tabular}{lrr}
\toprule
指标 & 正式23架次 & 22架次见证\\
\midrule
A/B/C架次 & 11/6/6 & 11/4/7\\
加权迟到/s & 0 & 0\\
makespan/s & 5693.231489 & 7767.058323\\
总能耗/kWh & 66.214460 & 66.323024\\
最紧硬截止余量/s & 202.096058 & 645.017134\\
最低返航SOC/\% & 20.097383 & 20.097383\\
\bottomrule
\end{tabular}
\end{table}

22架次方案只证明“22架次可行”，不证明22架次的最小makespan，也不能据此推出所有22架次方案均慢于23架次主方案。

\subsection{V7全局有效上下界}

\begin{table}[H]
\centering
\caption{Q2-EVIDENCE-V7-FOCUSED认证摘要}
\label{tab:app_q2_v7}
\begin{tabular}{lr}
\toprule
证据项 & 数值\\
\midrule
全局有效下界/s & 5433.956428\\
认证可行上界/s & 5693.231490\\
区间宽度/s & 259.275062\\
UB归一化间隙 & 4.554093092\%\\
互补分支数 & 809\\
叶证书数 & 810\\
完整机型定价调用 & 2430\\
审计失败数 & 0\\
原问题makespan全局最优 & 未证明\\
\bottomrule
\end{tabular}
\end{table}



\subsection{有限邻域与计算层负结果}

围绕当前主方案共检查330个定向局部邻域，329个有限模型返回不可行，1个达到时限后保持未知；不能将该结果解释为原问题局部或全局最优证明。pricing dominance实验只评价固定真实定价输入：24个输入的各组中位耗时之和由6.323099 s降至2.026581 s，比值为3.1201；生成标签由5,706,982降至2,061,175，最小约化成本保持一致。该3.1201倍不是整个Q2求解器速度比。

% ============================================================
\section{问题三运输--中继联合方案与验证}
\label{app:q3}

\subsection{当前正式release摘要}

当前正式版本为 \texttt{Q3-BOTTLENECK-V6/time}，物理模型版本为 \texttt{Q3-HOVER-COMPATIBLE-V5}。

\begin{table}[H]
\centering
\caption{Q3 V6正式主方案摘要}
\label{tab:app_q3_summary}
\begin{tabular}{lr}
\toprule
指标 & 数值\\
\midrule
运输架次 & 23\\
中继架次 & 4\\
实体运输无人机 & 8\\
运输电池 & 14\\
中继无人机 & 2\\
中继能源组件 & 3\\
联合最后返航/s & 5836.969929\\
运输能耗/kWh & 66.251280\\
中继能耗/kWh & 2.688136\\
总能耗/kWh & 68.939416\\
加权迟到 & 0\\
硬迟到箱数 & 0\\
最小硬截止余量/s & 360.866938\\
最低运输返航SOC/\% & 20.097383\\
\bottomrule
\end{tabular}
\end{table}

\subsection{中继转场高度口径}

中继转场采用
\[
H_R(p)=\max\left\{\max_{c\in\mathcal C(O01,p)}Z(c)+50,\ z_p,\ z_{O01}\right\}.
\]
该规则只作用于中继往返转场，普通运输航段仍使用沿线最高地形+50 m。旧release中将悬停高度额外截断在地形净空巡航面以下，因此其有限候选域证书不能自动继承到V6。

\subsection{连续通信验证}

V6主方案重新构造真实运输轨迹、完整中继需求区间与临界端点。主验证器完成17,756项物理、资源、区间和端点检查，失败0；独立时空点核验36,672点，未发现断链。保存证据包含700份完整区间证书和1,238份临界端点证书。

\subsection{关键保障切换}

关键C型运输任务在V5中从相对956.973974 s起必须依赖西侧中继；V6将此前段保障延长至1039.150288 s，同时使西侧中继提前5.821245 s就绪。两项共同解释约87.997559 s的同模型工期改善。该分析解释保存日程中的关键链，不构成原问题全局下界。

\subsection{时间--能耗与鲁棒备选}

\begin{table}[H]
\centering
\caption{Q3 V6已验证候选}
\label{tab:app_q3_robust}
\begin{tabular}{lrrr}
\toprule
方案 & 额外统一损耗/dB & 联合工期/s & 总能耗/kWh\\
\midrule
time & 0 & 5836.969929 & 68.939416\\
energy & 0 & 5911.548135 & 68.735791\\
robust025 & 0.25 & 5866.140527 & 69.031751\\
q4\_tradeoff & 0 & 5866.140527 & 68.975497\\
\bottomrule
\end{tabular}
\end{table}

旧+0.50/+0.55 dB结果属于历史release，不属于V6；当前只把+0.25 dB方案作为V6下完整核验的通信裕度备选。

% ============================================================
\section{问题四V6严格分区与资源向量}
\label{app:q4}

\subsection{严格模型不可拆块}

冻结Q3 V6/time任务、绝对时刻和实际通信保障关系后，严格模型形成
\[
B_1=\{S001\},\qquad
B_2=\{S002,S003,S004,S005,S007,S008,S009,S010,S011,S012,S013,S014,S015\},\qquad
B_3=\{S006\}.
\]
严格两组共有3种无标签分区，严格三组只有1种。

\subsection{当前正式资源向量}

\begin{table}[H]
\centering
\caption{Q4 V6严格分类型资源需求与库存}
\label{tab:app_q4_vectors}
\begin{tabular}{lrrrr}
\toprule
资源类型 & 原库存 & 全局共享 & 严格K2 & 严格K3\\
\midrule
A型运输机 & 4 & 4 & 4 & 5\\
B型运输机 & 2 & 2 & 2 & 2\\
C型运输机 & 2 & 2 & 3 & 5\\
A型电池 & 6 & 6 & 6 & 7\\
B型电池 & 4 & 4 & 4 & 4\\
C型电池 & 4 & 4 & 4 & 6\\
中继无人机 & 2 & 2 & 2 & 2\\
中继能源组件 & 6 & 3 & 3 & 3\\
\midrule
总配置 & 30 & 27 & 28 & 34\\
\bottomrule
\end{tabular}
\end{table}

严格两组的正式需求和缺口为
\[
\mathbf D^{(2)}=(4,2,3,6,4,4,2,3),\qquad
\boxed{\mathbf G^{(2)}=(0,0,1,0,0,0,0,0)}.
\]
严格三组的需求和缺口为
\[
\mathbf D^{(3)}=(5,2,5,7,4,6,2,3),\qquad
\boxed{\mathbf G^{(3)}=(1,0,3,1,0,2,0,0)}.
\]

两组正式分区为“除S006外14区 / S006”，总缺口1；三组唯一分区为“S001 / 其余13区 / S006”，总缺口7。

\subsection{跨问折中}

V6另保留一个 \texttt{q4\_tradeoff} 候选：联合工期5866.140527 s、总能耗68.975497 kWh，其严格三组缺口为6。它在Q3时间--能耗二维上不优于正式time方案，但说明Q3最短工期与Q4最少独立配置并不等价。

\subsection{历史敏感性边界}

旧Q4-SENSITIVITY-V2中的relay-copy、N-1、Pareto和库存敏感性使用旧Q3日程，未针对V6重新计算。它们仍保留在仓库历史目录，但不能与当前V6 strict资源向量混用。

\subsection{资源最小数量的交叉核验}

对固定任务组和固定资源类型，最小资源数量
\[
N_{gp}^{\min}=\max_t n_{gp}(t)
\]
由事件扫描和半开区间峰值计算得到，并在V6完整运行包中以实体资源链和独立最大匹配进行交叉核验。

% ============================================================

\section{复现、版本锁与证据清单}
\label{app:repro}

\subsection{正式版本入口}

\begin{table}[H]
\centering
\caption{四问正式结果入口}
\label{tab:app_release}
\begin{tabular}{lll}
\toprule
问题 & 正式版本/提交 & 论文复现入口\\
\midrule
Q1 & \texttt{f905acc8482966c...} & \path{problems/D/q1/}\\
Q2 & Q2-EVIDENCE-V7-FOCUSED & \path{problems/D/q2/evidence_v7/}\\
Q3 & Q3-BOTTLENECK-V6 / time & \path{problems/D/q3/release_v6/time/}\\
Q4 & Q4-STRICT-FROM-Q3-V6 & \path{problems/D/q4/results_q3_v6/}\\
\bottomrule
\end{tabular}
\end{table}

\subsection{哈希锁定的伴随归档}

\begin{longtable}{p{6.5cm}p{8.0cm}}
\caption{正式伴随归档SHA-256}\label{tab:app_archives}\\
\toprule
归档 & SHA-256\\
\midrule
\endfirsthead
\toprule
归档 & SHA-256\\
\midrule
\endhead
D题\_Q2定向证据强化V7\_源码证书与复核.zip &
\texttt{0804f989699dc506b3111dacc7408d9acdeb6e1115e7f32ec36934116de6ce58}\\
D题\_Q3正式封口\_FINAL\_源码结果与交接.zip &
\texttt{97f5ea8018a557d604bee87d6886675ecf7648f035e8df0f47f20ff6e7fbc56b}\\
D题\_Q4稳健性完善V2\_源码结果与验证.zip &
\texttt{9c8d1a6410c2de0a410318fc3525dc162613deabf4ba9dd1d59cdeb98437f5b4}\\
D题\_Q2外部方案复核\_代码明细与验证.zip &
\texttt{b41a0416313a123e734e01d7ca1672b99acee41f16edd79812302e1007610b77}\\
\bottomrule
\end{longtable}

\subsection{证据等级与最终提交前检查}

正式论文中只允许使用与证据强度相匹配的表述：
\begin{itemize}
  \item Q1最少18架次：原Q1模型下全局证明；
  \item Q2加权迟到0：第一层目标全局证明；
  \item Q2 makespan：有效全局下界与完整可行上界，未闭合；
  \item Q3 V6：连续区间/临界端点与独立重放证明保存方案完整可行；
  \item Q3固定结构：连续时间LP证明当前固定离散结构下的时序最优；
  \item Q3鲁棒性：仅$+0.25$ dB为当前V6完整可行见证，旧$+0.50/+0.55$ dB不继承；
  \item Q4 K2/K3：冻结Q3 V6/time日程后的strict完整分区枚举。
\end{itemize}

最终提交前应再次以 \path{paper_integration/FINAL_SELECTION.json} 为唯一选择入口，逐项核对正文、摘要、图表、附录和机器可读结果，避免旧版本数字残留。
